% ============================================================================
% III. METHODOLOGY (REVISED - FINAL)
% ============================================================================

\section{Methodology}
\label{sec:methodology}

% TikZ pipeline diagram
\begin{figure*}[!t]
\centering
\begin{tikzpicture}[
    node distance=1.5cm,
    block/.style={rectangle, draw=black, thick, fill=blue!10, text width=2.2cm, minimum height=1.2cm, align=center, rounded corners=3pt},
    data/.style={rectangle, draw=black, thick, fill=green!10, text width=2cm, minimum height=1cm, align=center},
    process/.style={rectangle, draw=black, thick, fill=orange!10, text width=2.2cm, minimum height=1.2cm, align=center, rounded corners=3pt},
    output/.style={rectangle, draw=black, thick, fill=purple!10, text width=2cm, minimum height=1cm, align=center, rounded corners=3pt},
    arrow/.style={-{Stealth[length=3mm]}, thick}
]

% Input
\node[data] (input) {4D ME-HSI\\$(x,y,\lambdaem,\lambdaex)$};

% Preprocessing
\node[block, right=of input] (preprocess) {Preprocessing\\Normalization\\Rayleigh Cutoff};

% Autoencoder
\node[block, right=of preprocess] (encoder) {Parallel\\Encoder\\Branches};
\node[block, right=of encoder] (latent) {Shared\\Latent Space\\$\vect{z} \in \R^{20}$};
\node[block, right=of latent] (decoder) {Parallel\\Decoder\\Branches};

% Perturbation Analysis
\node[process, below=1.2cm of latent] (perturb) {Latent Space\\Perturbation\\Analysis};

% Influence scores
\node[process, left=of perturb] (influence) {Wavelength\\Influence\\Scores};

% MMR Selection
\node[process, left=of influence] (mmr) {MMR\\Diversity\\Selection};

% Output
\node[output, left=of mmr] (output) {Selected\\Wavelengths};

% Arrows
\draw[arrow] (input) -- (preprocess);
\draw[arrow] (preprocess) -- (encoder);
\draw[arrow] (encoder) -- (latent);
\draw[arrow] (latent) -- (decoder);
\draw[arrow] (latent) -- (perturb);
\draw[arrow] (perturb) -- (influence);
\draw[arrow] (influence) -- (mmr);
\draw[arrow] (mmr) -- (output);

% Labels
\node[above=0.3cm of encoder, font=\footnotesize\itshape] {Stage 1: Representation Learning};
\node[below=0.3cm of influence, font=\footnotesize\itshape] {Stage 2: Attribution};
\node[below=0.3cm of mmr, font=\footnotesize\itshape] {Stage 3: Selection};

% Brace for autoencoder
\draw[decorate, decoration={brace, amplitude=8pt, mirror, raise=4pt}]
    (encoder.south west) -- (decoder.south east) node[midway, below=14pt, font=\footnotesize] {3D CAE (Training)};

\end{tikzpicture}
\caption{Overview of the proposed wavelength selection framework. The pipeline comprises three stages: (1) a 3D CAE learns a compressed representation of 4D ME-HSI data; (2) perturbation analysis attributes latent dimensions to individual wavelengths; (3) MMR selection identifies diverse, informative band subsets.}
\label{fig:pipeline}
\end{figure*}

% ----------------------------------------------------------------------------
\subsection{Data Preprocessing}
\label{sec:preprocessing}

Raw ME-HSI acquisitions required several preprocessing steps before autoencoder training. (Figure~\ref{fig:pipeline})

\subsubsection{Intensity Normalization}
Fluorescence intensity depends on excitation source power and exposure time, which may vary across acquisitions.
Since the LED-based multi-excitation source operated at different power levels and exposure times for each excitation wavelength, raw intensities were normalized to ensure comparability across excitation channels:
\begin{equation}
I_{\text{norm}}(x, y, \lambdaem, \lambdaex) = \frac{I_{\text{raw}}(x, y, \lambdaem, \lambdaex)}{t_{\text{exp}}} \times \frac{t_{\text{ref}} \cdot P_{\text{ref}}}{P(\lambdaex)}
\label{eq:normalization}
\end{equation}
where $t_{\text{exp}}$ is the exposure time and $P(\lambdaex)$ is the lamp power at excitation wavelength $\lambdaex$.
The reference values $t_{\text{ref}}$ and $P_{\text{ref}}$ were set to the maximum exposure time and maximum lamp power, respectively, across all excitation wavelengths, ensuring that the normalization factor scales intensities upward for weaker or shorter-exposed channels.
Global min-max normalization was then applied to scale all values to $[0, 1]$.

\subsubsection{Rayleigh Scattering Removal}
Elastic scattering produces artifacts at two spectral locations.
First-order Rayleigh scattering occurs when emission wavelength approaches excitation wavelength ($\lambdaem \approx \lambdaex$), therefore, a spectral cutoff with offset $\Delta\lambda = 40$\, nm excluded emission bands where $\lambdaem < \lambdaex + \Delta\lambda$.
Second-order Rayleigh scattering produces artifacts near $\lambdaem \approx 2\lambdaex$, therefore emission bands within $\pm 40$\, nm of this harmonic were also excluded.
Together, these cutoffs resulted in variable numbers of valid emission bands per excitation, yielding 192 valid excitation-emission pairs from the original acquisition grid.
% \begin{figure}[!t]
% \centering
% \begin{tikzpicture}
% \begin{axis}[
%     width=0.48\textwidth,
%     height=0.38\textwidth,
%     xlabel={Emission Wavelength (nm)},
%     ylabel={Excitation Wavelength (nm)},
%     xmin=400, xmax=750,
%     ymin=305, ymax=435,
%     ytick={310, 325, 340, 355, 370, 385, 400, 415, 430},
%     xtick={400, 450, 500, 550, 600, 650, 700, 750},
%     tick label style={font=\footnotesize},
%     label style={font=\small},
%     grid=major,
%     grid style={dashed, gray!30},
%     clip=true,
% ]

% % 1st order exclusion zone (Em < Ex + 40): upper-left triangle
% \addplot[fill=red!25, draw=none] coordinates {
%     (400, 360) (475, 435) (400, 435)
% } -- cycle;

% % 1st order cutoff line: Em = Ex + 40
% \addplot[red, thick, dashed, samples=2, domain=355:435] ({x + 40}, {x});

% % 2nd order exclusion band: |Em - 2*Ex| < 40, diagonal band
% \addplot[fill=blue!20, draw=none] coordinates {
%     (570, 305) (650, 305) (750, 355) (750, 395)
% } -- cycle;

% % 2nd order boundary lines
% \addplot[blue, thick, dotted, samples=2, domain=305:395] ({2*x - 40}, {x});
% \addplot[blue, thick, dotted, samples=2, domain=305:355] ({2*x + 40}, {x});

% % Labels parallel to cutoff lines (visual angles computed from axis aspect ratio)
% \node[red, font=\scriptsize, rotate=65, anchor=south] at (axis cs:435, 398) {1\textsuperscript{st} order cutoff};
% \node[blue, font=\scriptsize, rotate=47, anchor=south] at (axis cs:660, 335) {2\textsuperscript{nd} order cutoff};
% \node[font=\scriptsize] at (axis cs:550, 400) {Valid};

% % Excitation wavelength tick markers
% \foreach \ex in {310, 325, 340, 365, 385, 400, 415, 430} {
%     \addplot[only marks, mark=|, mark size=2pt, thick, black] coordinates {(400, \ex)};
% }

% \end{axis}
% \end{tikzpicture}
% \caption{Rayleigh scattering cutoffs in the excitation-emission matrix (EEM) space. The first-order exclusion zone (red shading, dashed line: $\lambdaem = \lambdaex + 40$\,nm) removes emission bands too close to the excitation wavelength. The second-order exclusion band (blue shading, dotted lines: $\lambdaem = 2\lambdaex \pm 40$\,nm) removes harmonic artifacts. The unshaded region between these cutoffs contains the 192 valid excitation-emission pairs used in this study.}
% \label{fig:rayleigh}
% \end{figure}

\subsubsection{Background Masking}
Non-sample regions (background) provide no discriminative information and may bias the autoencoder toward reconstructing uninformative patterns.
Binary masks $M(x, y) \in \{0, 1\}$ were generated through intensity thresholding or manual annotation, excluding background pixels from both training and evaluation.

% ----------------------------------------------------------------------------
\subsection{3D CAE}
\label{sec:architecture}

The central architectural challenge in processing 4D ME-HSI data was the variable spectral dimension across excitations.
After the Rayleigh cutoff, each excitation wavelength $\lambdaex$ yielded a different number of valid emission bands $N_{\text{em}}(\lambdaex)$.
A monolithic architecture would require padding and masking complexities to address the varying emission wavelength counts for each excitation; instead, a parallel-branch design was employed where each excitation is processed by a dedicated encoder branch, with features merged into a shared latent representation.

\subsubsection{Encoder Architecture}
For each excitation wavelength $\lambdaex^{(i)}$, $i = 1, \ldots, N_{\text{ex}}$, an encoder branch was defined, comprising a 3D convolutional layer that processes the spatial-spectral cube:
\begin{equation}
\vect{h}^{(i)} = \sigma\left(\text{Conv3D}\left(\vect{X}^{(i)}\right)\right)
\label{eq:encoder}
\end{equation}
where $\vect{X}^{(i)} \in \R^{H \times W \times N_{\text{em}}^{(i)}}$ is the input cube for excitation $i$, Conv3D applies $k_1 = 20$ filters of size $5 \times 5 \times \min(5, N_{\text{em}}^{(i)})$, and $\sigma$ is the sigmoid activation function.
The subscripts in $k_1$ and $k_3$ (introduced below) index the first and third convolutional layers; the intermediate feature-merging step has no learned parameters.
The adaptive spectral kernel size accommodated variable band counts across excitations.

After convolution, adaptive average pooling standardized the spectral dimension to 1, yielding feature maps $\vect{h}^{(i)} \in \R^{k_1 \times H \times W}$ with consistent dimensions across excitations.
Dropout regularization (rate 0.5) was applied after pooling to prevent overfitting.

\subsubsection{Feature Merging}
The excitation-specific feature maps were combined through element-wise averaging:
\begin{equation}
\bar{\vect{h}} = \frac{1}{N_{\text{ex}}} \sum_{i=1}^{N_{\text{ex}}} \vect{h}^{(i)}
\label{eq:merge}
\end{equation}
This merging strategy offers several advantages over alternatives.
Concatenation would yield feature dimensionality that grows with $N_{\text{ex}}$, increasing model complexity.
Averaging provided implicit noise reduction and was robust to intensity variations across excitations.
The merged features $\bar{\vect{h}}$ captured cross-excitation correlations by integrating patterns from all excitation branches.

\subsubsection{Shared Latent Space}
An additional convolutional layer transformed the merged features into the latent representation:
\begin{equation}
\vect{z} = \sigma\left(\text{Conv3D}\left(\bar{\vect{h}}\right)\right)
\label{eq:latent}
\end{equation}
where $\vect{z} \in \R^{k_3 \times 1 \times H \times W}$ with $k_3 = 20$ channels constitutes the bottleneck encoding.
This latent space $\vect{z}$ encapsulates the essential structure of the 4D ME-HSI data in a compressed form.
The perturbation analysis (Section~\ref{sec:perturbation}) probes this space to identify which wavelengths contribute most to the learned representation.

\subsubsection{Decoder Architecture}
The decoder mirrors the encoder with symmetric structure.
A shared decoding layer transformed the latent representation:
\begin{equation}
\vect{g} = \sigma\left(\text{Conv3D}\left(\vect{z}\right)\right)
\label{eq:decode_shared}
\end{equation}
yielding features $\vect{g} \in \R^{k_1 \times 1 \times H \times W}$.
Parallel decoder branches then reconstructed each excitation's spectral cube:
\begin{equation}
\hat{\vect{X}}^{(i)} = \sigma\left(\text{Conv3D}^{(i)}\left(\vect{g}\right)\right)
\label{eq:decode_branch}
\end{equation}
where $\text{Conv3D}^{(i)}$ maps from $k_1$ to $N_{\text{em}}^{(i)}$ channels.
Figure~\ref{fig:architecture} illustrates the complete architecture.

\begin{figure*}[!t]
\centering
\resizebox{\textwidth}{!}{%
\begin{tikzpicture}[
    scale=1,
    transform shape,
    block/.style={rectangle, draw, thick, fill=blue!10, minimum width=1.5cm, minimum height=0.6cm, align=center, font=\footnotesize},
    arrow/.style={-{Stealth[length=2mm]}, thick},
    label/.style={font=\scriptsize}
]

% Input branches
\node[block, fill=green!20] at (0, 2) (in1) {$\vect{X}^{(1)}$};
\node[block, fill=green!20] at (0, 1) (in2) {$\vect{X}^{(2)}$};
\node at (0, 0.3) {$\vdots$};
\node[block, fill=green!20] at (0, -0.5) (inN) {$\vect{X}^{(N)}$};

% Encoder branches
\node[block] at (2, 2) (enc1) {Conv3D};
\node[block] at (2, 1) (enc2) {Conv3D};
\node at (2, 0.3) {$\vdots$};
\node[block] at (2, -0.5) (encN) {Conv3D};

% Merge
\node[block, fill=orange!20] at (4, 0.75) (merge) {Average};

% Latent
\node[block, fill=red!20] at (6, 0.75) (latent) {$\vect{z}$};

% Shared decode
\node[block, fill=orange!20] at (8, 0.75) (dec_shared) {Conv3D};

% Decoder branches
\node[block] at (10, 2) (dec1) {Conv3D};
\node[block] at (10, 1) (dec2) {Conv3D};
\node at (10, 0.3) {$\vdots$};
\node[block] at (10, -0.5) (decN) {Conv3D};

% Output branches
\node[block, fill=green!20] at (12, 2) (out1) {$\hat{\vect{X}}^{(1)}$};
\node[block, fill=green!20] at (12, 1) (out2) {$\hat{\vect{X}}^{(2)}$};
\node at (12, 0.3) {$\vdots$};
\node[block, fill=green!20] at (12, -0.5) (outN) {$\hat{\vect{X}}^{(N)}$};

% Arrows
\foreach \i in {1,2,N} {
    \draw[arrow] (in\i) -- (enc\i);
    \draw[arrow] (dec\i) -- (out\i);
}
\draw[arrow] (enc1) -- (merge);
\draw[arrow] (enc2) -- (merge);
\draw[arrow] (encN) -- (merge);
\draw[arrow] (merge) -- (latent);
\draw[arrow] (latent) -- (dec_shared);
\draw[arrow] (dec_shared) -- (dec1);
\draw[arrow] (dec_shared) -- (dec2);
\draw[arrow] (dec_shared) -- (decN);

% Labels
\node[above=0.1cm of enc1, label] {Parallel Encoders};
\node[above=0.1cm of dec1, label] {Parallel Decoders};
\node[below=0.3cm of latent, label, text=red] {Bottleneck};

\end{tikzpicture}%
}
\caption{3D CAE architecture with parallel branches. Each excitation is processed by a dedicated encoder branch; features are averaged into a shared latent space; parallel decoder branches reconstruct the original 4D structure.}
\label{fig:architecture}
\end{figure*}

\subsubsection{Loss Function and Training}
The autoencoder was trained to minimize masked mean squared error:
\begin{equation}
\mathcal{L} = \frac{1}{N_{\text{valid}}} \sum_{i=1}^{N_{\text{ex}}} \sum_{x,y} M(x,y) \cdot \left\| \vect{X}^{(i)}_{x,y} - \hat{\vect{X}}^{(i)}_{x,y} \right\|^2
\label{eq:loss}
\end{equation}
where $M(x,y)$ is the background mask, $N_{\text{valid}}$ is the number of valid pixels, and the sum runs over all excitations and spatial locations.
Masking ensured the network focused on sample regions rather than background.

Training employed the Adam optimizer~\cite{kingma2014adam} with learning rate $\eta = 0.001$, learning rate scheduling (ReduceOnPlateau with factor 0.5 and patience 5), and early stopping (patience 5 epochs).
Large images were partitioned into $64 \times 64$ patches with 8-pixel overlap for batch training; typical convergence occurred within 25--30 epochs.

% ----------------------------------------------------------------------------
\subsection{Perturbation-Based Wavelength Attribution}
\label{sec:perturbation}

The trained autoencoder learned a compressed representation that captures the essential structure of ME-HSI data.
Different latent dimensions encode different aspects of this structure; the goal was to attribute each latent dimension's influence back to specific excitation-emission wavelength pairs.
If perturbing a latent dimension significantly changed the reconstruction of a particular wavelength, that wavelength must be encoded by that dimension.

\subsubsection{Latent Dimension Scoring}
Not all latent dimensions carry equal information.
The $k_3 = 20$ latent dimensions were first ranked by importance using variance-based scoring:
\begin{equation}
\text{score}(d) = \text{Var}_{\text{samples}}\left(\vect{z}[\cdot, d, \cdot, \cdot, \cdot]\right)
\label{eq:dim_score}
\end{equation}
where $\vect{z}[\cdot, d, \cdot, \cdot, \cdot]$ denotes dimension $d$ across all samples and spatial locations.
Higher variance indicated that the dimension captured meaningful variation rather than constant or near-zero values.
The top $k$ dimensions were then selected for perturbation analysis, where $k$ is treated as a hyperparameter; empirical evaluation (Section~\ref{sec:results}) revealed that $k=1$ (using only the highest-variance dimension) achieved optimal performance.

Alternative scoring methods include activation magnitude (mean absolute value) and reconstruction sensitivity (MSE when dimension is zeroed).
Variance-based scoring yielded the most consistent wavelength rankings, likely because high-variance dimensions captured the primary modes of spectral variation that distinguish different materials.

\subsubsection{Systematic Perturbation}
For each selected latent dimension $d$, the standard deviation $\sigma_d$ was computed across spatial locations and samples, and then the perturbed representations were generated:
\begin{align}
\vect{z}^{+}_d &= \vect{z} + \epsilon \cdot \sigma_d \cdot \vect{e}_d \label{eq:perturb_pos}\\
\vect{z}^{-}_d &= \vect{z} - \epsilon \cdot \sigma_d \cdot \vect{e}_d \label{eq:perturb_neg}
\end{align}
where $\vect{e}_d$ is a one-hot vector selecting dimension $d$, and $\epsilon$ controls perturbation magnitude.
Values of $\epsilon \in \{15, 30, 45\}$ were used, and results across scales were aggregated to ensure robustness.
The $\sigma$-based scaling ensured perturbations were proportional to the natural variation in each dimension, avoiding both negligibly small and destructively large perturbations.

\subsubsection{Sensitivity Mapping}
Decoding the perturbed representations yielded reconstructions $\hat{\vect{X}}^{+}_d$ and $\hat{\vect{X}}^{-}_d$.
The sensitivity of each wavelength to perturbation in dimension $d$ was computed:
\begin{equation}
\Delta_d(\lambdaex, \lambdaem) = \frac{\left| \hat{X}^{+}_d(\lambdaex, \lambdaem) - \hat{X}^{-}_d(\lambdaex, \lambdaem) \right|}{2\epsilon \cdot \sigma_d}
\label{eq:sensitivity}
\end{equation}
where averaging was performed over spatial locations and samples.
This finite difference approximation estimates the partial derivative of reconstruction with respect to latent dimension $d$, quantifying how strongly each wavelength's reconstruction responds to changes in that dimension.

\subsubsection{Influence Accumulation}
Sensitivity was accumulated across latent dimensions, weighted by dimension importance:
\begin{equation}
\text{Influence}(\lambdaex, \lambdaem) = \sum_{d \in \text{top-}k} \text{score}(d) \cdot \overline{\Delta_d(\lambdaex, \lambdaem)}
\label{eq:influence}
\end{equation}
where $\overline{\Delta_d}$ denotes the mean sensitivity across perturbation magnitudes.
The result is an influence matrix $\mat{I} \in \R^{N_{\text{ex}} \times \max(N_{\text{em}})}$ assigning importance scores to each excitation-emission pair.

% ----------------------------------------------------------------------------
\subsection{Maximum Marginal Relevance Selection}
\label{sec:mmr}

Selecting the top-$K$ wavelengths by influence score alone may yield redundant bands.
Adjacent wavelengths often captured correlated fluorescence patterns; a diverse selection covering different spectral regions may be more informative than a concentrated selection.
Maximum Marginal Relevance (MMR)~\cite{carbonell1998use} was adapted from information retrieval to balance relevance with diversity.

\subsubsection{MMR Formulation}
Given the set of all wavelength pairs $\mathcal{W}$ and a partial selection $\mathcal{S}$, MMR scores the next candidate as:
\begin{equation}
\text{MMR}(w_i) = \lambda \cdot \text{Rel}(w_i) - (1 - \lambda) \cdot \max_{w_j \in \mathcal{S}} \text{Sim}(w_i, w_j)
\label{eq:mmr}
\end{equation}
where $\text{Rel}(w_i)$ is the normalized influence score, $\text{Sim}(w_i, w_j)$ is the cosine similarity between spectral profiles at wavelengths $w_i$ and $w_j$, and $\lambda \in [0, 1]$ controls the relevance-diversity trade-off.

At $\lambda = 1$, MMR reduces to pure relevance ranking; at $\lambda = 0$, selection maximizes diversity without regard for informativeness.
Empirical evaluation (Section~\ref{sec:results}) identified $\lambda = 0.5$ as optimal, providing balanced consideration of both criteria.

\subsubsection{Selection Procedure}
Algorithm~\ref{alg:mmr} presents the greedy MMR selection procedure.
The algorithm iteratively selected the wavelength that maximized the MMR score until $K$ wavelengths were chosen.
The output is an ordered list of excitation-emission pairs, ranked by selection order, enabling flexible thresholding based on accuracy-efficiency trade-offs.

\begin{algorithm}[t]
\caption{MMR Wavelength Selection}
\label{alg:mmr}
\begin{algorithmic}[1]
\REQUIRE Wavelength set $\mathcal{W}$, influence matrix $\mat{I}$, target count $K$, trade-off parameter $\lambda$
\ENSURE Selected wavelength set $\mathcal{S}$
\STATE Initialize $\mathcal{S} \gets \emptyset$
\STATE Compute $\text{Rel}(w) \gets \mat{I}(w) / \max(\mat{I})$ for all $w \in \mathcal{W}$
\FOR{$i = 1$ to $K$}
    \FOR{each $w \in \mathcal{W} \setminus \mathcal{S}$}
        \IF{$\mathcal{S} = \emptyset$}
            \STATE $\text{MMR}(w) \gets \text{Rel}(w)$
        \ELSE
            \STATE $\text{MMR}(w) \gets \lambda \cdot \text{Rel}(w) - (1-\lambda) \cdot \max_{w_j \in \mathcal{S}} \text{Sim}(w, w_j)$
        \ENDIF
    \ENDFOR
    \STATE $w^* \gets \arg\max_{w \in \mathcal{W} \setminus \mathcal{S}} \text{MMR}(w)$
    \STATE $\mathcal{S} \gets \mathcal{S} \cup \{w^*\}$
\ENDFOR
\RETURN $\mathcal{S}$
\end{algorithmic}
\end{algorithm}

% ----------------------------------------------------------------------------
\subsection{Validation Protocol}
\label{sec:validation}

The wavelength selection pipeline (Sections~\ref{sec:architecture}--\ref{sec:mmr}) is entirely unsupervised and operates without access to class labels.
A supervised classification algorithm is used to quantify how much discriminative power remains in the compressed features; the classification described below serves \emph{only} to evaluate the quality of selected wavelengths---it does not influence the selection process.
In a deployment scenario, wavelength selection would proceed identically whether or not labeled data exists.

For each pixel $(x, y)$ with ground-truth labels, spectral values at the $K$ selected wavelength pairs form the feature vector:
\begin{equation}
\resizebox{0.88\columnwidth}{!}{$
\vect{f}(x, y) = \left[ I(x, y, \lambda_{\mathrm{em}}^{(1)}, \lambda_{\mathrm{ex}}^{(1)}), \ldots, I(x, y, \lambda_{\mathrm{em}}^{(K)}, \lambda_{\mathrm{ex}}^{(K)}) \right]^\top$},
\label{eq:features}
\end{equation}
where $(\lambdaem^{(k)}, \lambdaex^{(k)})$ is the $k$-th selected pair.
Features were standardized (zero mean, unit variance) and classified using a $k$-nearest neighbors (KNN) classifier~\cite{cover1967nearest} with 5 neighbors and Euclidean distance.
This non-parametric classifier was chosen to isolate the effect of wavelength selection from classifier capacity.
A baseline using all 192 bands with the same classifier established the performance reference.
Cohen's Kappa ($\kappa$)~\cite{cohen1960coefficient}, classification accuracy, and weighted F1-score served as primary evaluation metrics; the complete metric set is defined in Section~\ref{sec:metrics}.